## TEMP holding note — RQ3 GNNWR pooled results (2026-08-11)

**Cluster fit+evaluate ran**: 2026-08-10/11 (`step2_rq3_gnnwr.sh`, 30 jobs — 3 sets x 2 cohorts x
5 folds, VIF-screened `nested_set2_top10`/`nested_set3_gated_terrain_wind_vif`/
`nested_set4_gated_all_vif`). All 30/30 completed, zero `FAILED` entries. Synced via
`step2_sync_results.sh`. **Pooling script never existed for the new tiers until today** — a
pre-existing script (`models/growth_curve_attribution/pool_gnnwr_kfold_results.py`, originally
built for the old named `SCOPES` system) turned out to already work unmodified for the new tiers,
since `--scope` is only ever used as a bare filename-pattern string, never validated against
`SCOPES` internally. Only the CLI's `choices=` list needed widening (additive, same pattern as
`run_baselines_env.py`'s earlier fix) — no logic changes.

### How these numbers were generated

- GNNWR fit+evaluate are combined in one cluster call per fold (the model is too large to sync
  down and reload locally, so evaluation happens on the cluster by design) — only each fold's
  `*_test_predictions.csv` comes back.
- `pool_gnnwr_kfold_results.py`: concatenates all 5 folds' test predictions per (set, cohort),
  computes ONE pooled R2 over the whole population (same "every plot's out-of-fold prediction,
  all folds concatenated" convention as `spatial_cv_check.py`'s EN/XGBoost pooling and RQ1's
  `kfold_summary`), plus per-fold R2 mean/SD for the precision comparison.
- Command: `python -m models.growth_curve_attribution.pool_gnnwr_kfold_results --cohort <cohort> --scope <set_name> --reference-set-size 0 --k-folds 5`
  (`--reference-set-size 0` = full population, matching how these jobs were actually run).

### Results — pooled + per-fold R2, all 6 (set x cohort) combinations

| Set | Cohort | n plots | n compartments | Pooled R2 | Per-fold R2 (mean±SD) |
|---|---|---:|---:|---:|---|
| Set2 | 4survey | 56,135 | 231 | 0.320 | 0.307±0.084 |
| Set2 | 6survey | 13,335 | 47 | -0.044 | -0.074±0.174 |
| Set3 | 4survey | 56,526 | 231 | 0.319 | 0.306±0.065 |
| Set3 | 6survey | 13,335 | 47 | 0.058 | 0.013±0.082 |
| Set4 | 4survey | 56,114 | 231 | 0.294 | 0.277±0.076 |
| Set4 | 6survey | 13,335 | 47 | -0.075 | -0.113±0.224 |

### Headline finding — the actual RQ3 question, now answerable

**Does GNNWR's spatially-varying fit beat global EN/XGBoost?** Comparing against
`TEMP_rq3_en_xgb_results_2026-08-11.tex`'s pooled R2:

| Set | EN (4survey) | XGB (4survey) | GNNWR (4survey) | EN (6survey) | XGB (6survey) | GNNWR (6survey) |
|---|---:|---:|---:|---:|---:|---:|
| Set2 | 0.286 | 0.266 | **0.320** | 0.057 | 0.031 | -0.044 |
| Set3 | 0.274 | 0.281 | **0.319** | 0.031 | 0.086 | 0.058 |
| Set4 | 0.240 | 0.250 | **0.294** | 0.056 | 0.015 | -0.075 |

**Yes on 4survey. On 6survey — see the reseed below, this single-seed number is not reliable
enough to call it a loss.** GNNWR beats both global models on EVERY set for 4survey (by 0.03-0.05
R2, a real and consistent margin, not noise-level given the fold SDs are ~0.07-0.08). The 6survey
numbers in this table (from seed 42 only) looked like a clean loss at first, but the 2026-08-12
reseed below shows the sign flips across seeds for every set — treat 6survey as "underpowered,
can't draw a directional conclusion" rather than "GNNWR loses here," and read the 4survey result
(unaffected by this) as the reliable part of this comparison.

### 6survey reseed (2026-08-12) — CORRECTS the "GNNWR loses on 6survey" framing above

Ran 2 more seeds (43, 44) on 6survey only (4survey's own per-fold SD already looked stable enough
not to need this). **Caught and fixed a real bug first**: `run_gnnwr()`'s `run_name` never
included `split_seed` at all — reseeding would have silently overwritten the seed-42 results with
an identical filename before ever syncing. Fixed additively (only appends `_seed{N}` when it
differs from the project default 42, verified byte-identical against the actual already-synced
seed-42 filename before running anything). Same fix applied to
`pool_gnnwr_kfold_results.py`, which had the identical gap.

| Set | Seed 42 | Seed 43 | Seed 44 | Mean | SD |
|---|---:|---:|---:|---:|---:|
| Set2 | -0.044 | 0.059 | 0.028 | 0.014 | 0.053 |
| Set3 | 0.058 | -0.002 | 0.108 | 0.054 | 0.055 |
| Set4 | -0.075 | 0.055 | 0.027 | 0.002 | 0.068 |

**The sign flips across seeds for every single set.** This is not "GNNWR reliably loses on
6survey" — it's "GNNWR's pooled R2 on 6survey is indistinguishable from zero and highly unstable
seed-to-seed," a meaningfully different and more accurate finding. The across-seed mean for every
set sits close to 0 (0.002–0.054), well within the seed-to-seed SD (0.05–0.07) — i.e. by the same
"is the gap bigger than the SD" test used elsewhere in this project (RQ1's winner check), **none
of these 6survey numbers clear their own noise floor**. Read this as "GNNWR doesn't work well
enough to draw a conclusion from on 6survey" rather than "GNNWR is worse than EN/XGBoost on
6survey" — the earlier framing overstated a directional claim that a single seed can't actually
support. 4survey's own headline (GNNWR beats EN/XGBoost on every set) is unaffected by this —
that result already had tight, seed-42-only fold SDs (0.065-0.084) an order of magnitude smaller
relative to its effect size than anything seen on 6survey here.

### Bootstrap 95% CIs on 4survey's pooled R2 (2026-08-15) — no retraining needed

Cluster-bootstrap (`models/common/bootstrap_ci.py`, unchanged, `n_resamples=1000` matching RQ1/
RQ2b's own default) computed directly from the already-saved `*_test_predictions.csv` files (which
already carry `cpmt`) — GNNWR's own training cost is irrelevant here, only the already-fitted
predictions are resampled.

| Set (4survey) | GNNWR R2 [95% CI] |
|---|---|
| Set2 | 0.320 [0.270, 0.369] |
| Set3 | 0.319 [0.263, 0.365] |
| Set4 | 0.294 [0.236, 0.347] |

### EN/XGBoost's own bootstrap 95% CI (2026-08-15) — closes the gap flagged above

Same method (`cluster_bootstrap_ci()`, unchanged, `n_resamples=1000`), applied to the already-saved
`predictions.csv` for RQ3's EN/XGBoost (Set2/3/4, 4survey) — checks whether GNNWR's win survives
CI-overlap scrutiny the way RQ1's winner-pick did, rather than resting on point estimates alone.

| Set | EN R2 [95% CI] | XGB R2 [95% CI] | GNNWR R2 [95% CI] (from above) |
|---|---|---|---|
| Set2 | 0.286 [0.241, 0.325] | 0.266 [0.214, 0.308] | 0.320 [0.270, 0.369] |
| Set3 | 0.274 [0.223, 0.316] | 0.281 [0.235, 0.320] | 0.319 [0.263, 0.365] |
| Set4 | 0.240 [0.192, 0.286] | 0.250 [0.201, 0.295] | 0.294 [0.236, 0.347] |

**GNNWR's CI overlaps EN's and XGBoost's CI in every single set** — by the same non-overlapping-CI
bar RQ1 used to pick its own winner, this specific test does not clear it. Point estimates are
consistently higher for GNNWR (every set, by 0.03-0.05), and this is independently corroborated by
the residual Moran's I check below (GNNWR reduces spatial clustering EN/XGBoost leave behind on
4survey) — so the honest read is "a real, plausible, corroborated edge that doesn't clear the
single strictest bar this project applies elsewhere," not a clean, CI-confirmed win the way RQ1's
was. The primer's item 1 (`a_draft_plan_14th_aug.md`) has been updated to reflect this, softening
"the one reliable positive result" language accordingly.

### Residual Moran's I vs. EN/XGBoost (2026-08-15) — current tier, closes the gap flagged in RQ2b

Same method as the RQ2b check (`compute_residual_morans_i()`, k=8, 999 permutations, seed=42,
`models/growth_curve_attribution/residual_spatial_autocorrelation_check.py`, unchanged), run for
real on RQ3's own current-tier EN/XGBoost and GNNWR residuals (`local_y_max_difference` target),
all 3 sets, both cohorts. No new fitting — all from already-saved `predictions.csv`/
`test_predictions.csv` files.

| Set / cohort | EN | XGB | GNNWR | GNNWR vs. best global |
|---|---:|---:|---:|---:|
| Set2 / 4survey | 0.7032 | 0.7059 | 0.6829 | -0.0203 |
| Set2 / 6survey | 0.5964 | 0.6064 | 0.6354 | +0.0390 |
| Set3 / 4survey | 0.7069 | 0.6958 | 0.6843 | -0.0115 |
| Set3 / 6survey | 0.6060 | 0.5721 | 0.6031 | +0.0310 |
| Set4 / 4survey | 0.7179 | 0.7085 | 0.6962 | -0.0123 |
| Set4 / 6survey | 0.5996 | 0.6091 | 0.6517 | +0.0521 |

(All p=0.001, the permutation floor, in every cell — clustering is real and significant
throughout, the question here is only the magnitude/direction of GNNWR's effect on it.)

**A clean, consistent, two-directional pattern — independent confirmation of two things already
established by other means, via a completely different diagnostic (spatial residual clustering,
not R2):**

1. **On 4survey, GNNWR reduces residual clustering relative to the best global model in all 3
   sets** (-0.012 to -0.020, small but consistent in direction every time). This corroborates the
   4survey accuracy win (GNNWR beats EN/XGBoost pooled R2) with independent evidence: it isn't just
   fitting the data better on average, it's specifically capturing spatial structure the global
   models leave behind — exactly the mechanism the accuracy win would imply if it's real.
2. **On 6survey, GNNWR *increases* residual clustering relative to the best global model in all 3
   sets** (+0.031 to +0.052, consistently the wrong direction). This independently corroborates the
   6survey reseed finding (R2 sign-flips across seeds, unreliable) via a completely different
   diagnostic — not just "the accuracy number is unstable," but "the spatial model appears to be
   fitting noise rather than real local structure" on the smaller cohort, consistent with GNNWR's
   spatial-weighting mechanism needing more reference density (47 vs. 231 compartments) than
   6survey provides.

**This closes the bridge flagged in RQ2b's own Moran's I finding** (`TEMP_rq2_attribution`,
current-tier EN/XGBoost residual Moran's I on `mean_cr_residual`, 0.65-0.74, barely moved by more
environmental data): RQ3's GNNWR modestly but consistently reduces that same kind of unexplained
spatial structure on 4survey, where the comparison is reliable — the motivation now has an
answered, not just hypothesised, other half.

### Not yet done

Local-coefficient maps (per-row `coef_*` columns already in each `test_predictions.csv`, not yet
plotted), Set4x6survey's known 17-vs-20-column parity gap (flagged 2026-08-11, may partly explain
why Set4 is GNNWR's worst 6survey result), cross-referencing GNNWR's local coefficients against
EN's global coefficients/XGBoost's SHAP values for the same variables, 4survey reseed (not done,
judged unnecessary given already-tight fold SDs), bootstrap CI on EN/XGBoost's own RQ3 numbers for
a direct CI-vs-CI comparison against GNNWR (see section above).
